\makeatletter
\def\input@path{{./}{../}{../../}{../../../}{../../../../}{preamble/}{../preamble/}{../../preamble/}{../../../preamble/}{../../../../book/preamble/}}
\makeatother
\input{preamble_exams.tex}
\input{graphs/highlights.tex}
% \GTenHighlightsfalse % Uncomment to disable instructional highlights.
% Population/model, sample evidence, design threshold, and result roles.
\newcommand{\Pop}[1]{\HighlightStroke{GTenHLParameter}{#1}}
\newcommand{\Obs}[1]{\HighlightStroke{GTenHLValue}{#1}}
\newcommand{\Design}[1]{\HighlightStroke{GTenHLThird}{#1}}
\newcommand{\Result}[1]{\HighlightStroke{GTenHLOutcome}{#1}}

\title{Grade 11 Global Economics\\Term 3 --- Lesson 1\\Home Practice C2: Hypotheses and tail selection}
\author{}\date{}
\hypersetup{pdftitle={Grade 11 Global Economics - Term 3 - Lesson 1 - C2},pdfsubject={Hypotheses and tail selection}}
\begin{document}
\maketitle
\setcounter{tocdepth}{1}
\setcounter{secnumdepth}{2}
\phantomsection\label{toc}
\tableofcontents
\section*{Evaluated criteria}
\begin{description}
\item[\textbf{C2}] Formulates null and alternative hypotheses for a population parameter and determines whether the corresponding test is left-tailed, right-tailed, or two-tailed from the direction of the research claim.
\end{description}
\textbf{How to practise.} Underline the research claim before considering sample results. Attempt each task before opening its solution link. Use the return arrows to revisit the problem, group, or contents. All datasets and scenarios in this activity are illustrative.
\clearpage
\section{Introduction}\label{sec:introduction}
\subsection*{Highlight legend}
\phantomsection\label{tab:highlight-legend}
\begin{tabularx}{\linewidth}{@{}>{\raggedright\arraybackslash}p{.22\linewidth}>{\raggedright\arraybackslash}p{.25\linewidth}>{\raggedright\arraybackslash}X@{}}
\toprule
\textbf{Symbol or notation}&\textbf{Concept}&\textbf{Brief meaning}\\\midrule
$\Pop{\mu},\Pop{p}$ & \Pop{Parameter} & Population feature. \\
$\Pop{H_0}$ & \Pop{Null} & Benchmark model. \\
$\Result{H_a}$ & \Result{Alternative} & Research direction. \\
$\Obs{\bar x},\Obs{\hat p}$ & \Obs{Evidence} & Observed summaries. \\
\bottomrule
\end{tabularx}
\subsection{Introductory example 1: Lower transfer fees}\label{ex:intro-one}
Before collecting data, a cooperative asks whether population mean fee is below 8 USD: $H_0:\mu=8$, $H_a:\mu<8$, left-tailed. A later $\bar x=8.2$ USD does not change that direction.
\subsection{Introductory example 2: Changed participation}\label{ex:intro-two}
An economist asks whether population participation differs from 60\%: $H_0:p=0.60$, $H_a:p\ne0.60$, two-tailed. A later $\hat p=0.64$ still belongs to that investigation.
\section{Notation Reference and Practice Procedure}\label{sec:notation-and-procedure}
\subsection{Notation reference}\label{tab:notation-reference}
\begin{tabularx}{\linewidth}{@{}>{\raggedright\arraybackslash}p{.16\linewidth}>{\raggedright\arraybackslash}p{.18\linewidth}>{\raggedright\arraybackslash}p{.28\linewidth}>{\raggedright\arraybackslash}X@{}}
\toprule
\textbf{Symbol or notation}&\textbf{Concept}&\textbf{Definition}&\textbf{Meaning in the introductory examples}\\\midrule
$\Pop{\mu},\Pop{p}$ & \Pop{Parameter} & A population mean or proportion, distinct from a sample estimate. & \textbullet\ \hyperref[ex:intro-one]{Example 1}: mean fee.\par\textbullet\ \hyperref[ex:intro-two]{Example 2}: participation rate. \\
$\Pop{H_0}$ & \Pop{Null} & Equality is the calculation boundary for a one-sided composite null. & \textbullet\ \hyperref[ex:intro-one]{Example 1}: $\mu=8$.\par\textbullet\ \hyperref[ex:intro-two]{Example 2}: $p=0.60$. \\
$\Result{H_a}$ & \Result{Alternative} & $<$ means left; $>$ means right; $\ne$ means two tails. & \textbullet\ \hyperref[ex:intro-one]{Example 1}: $\mu<8$.\par\textbullet\ \hyperref[ex:intro-two]{Example 2}: $p\ne0.60$. \\
$\Obs{\bar x},\Obs{\hat p}$ & \Obs{Evidence} & Sample results inform decisions, not the preselected research direction. & \textbullet\ \hyperref[ex:intro-one]{Example 1}: 8.2 USD.\par\textbullet\ \hyperref[ex:intro-two]{Example 2}: 0.64. \\
\bottomrule
\end{tabularx}
\subsection{Practice procedure}
\begin{enumerate}
\item Define the population parameter, benchmark, units, and period.
\item Read the research claim before looking at sample results.
\item Write the equality for $H_0$ and $<$, $>$, or $\ne$ for $H_a$ from the claim.
\item Name the tail and explain the claim words.
\item Keep direction fixed. Explain a composite null for a one-sided guarantee.
\end{enumerate}
\section{General Worked Problems}\label{sec:general-worked}
\SectionProblemsTableOfContents{\SectionProblemLink{prob:claims}{Three participation questions} \\
\SectionProblemLink{prob:notation}{Costs and returns}}
\subsection{Three participation questions}\label{prob:claims}
The online-trading benchmark is 62\% of eligible firms. Separate preregistered studies ask whether participation is (a) lower, (b) higher, or (c) different. Each uses independent random sampling and $\alpha=0.05$. A pilot found 112 users among 200 firms.
\QuestionsTasks[sec:general-worked]{{2}/{sol:claims}/{Write all hypothesis pairs and tails; explain why the pilot does not select direction.}}
\subsubsection*{Worked solution}\label{sol:claims}
Let $\Pop{p}$ be population online-trading proportion. Each calculation boundary is $H_0:p=0.62$. In (a), $\Result{H_a}:p<0.62$, left-tailed, because the claim is lower. In (b), $H_a:p>0.62$, right-tailed, because it is higher. In (c), $H_a:p\ne0.62$, two-tailed, because either direction matters. Composite one-sided nulls are $p\ge0.62$ and $p\le0.62$. The pilot's $\Obs{\hat p}=0.56$ is sample evidence. Choosing a left tail solely from this observation changes the question after examining evidence. No test decision is requested.
\secInicio[prob:claims][sec:general-worked]
\subsection{Costs and returns}\label{prob:notation}
A lender asks whether mean processing cost exceeds 75 USD. Another analyst asks whether population mean annual return differs from 4\%, measured in percentage points. A draft says $H_0:\bar x=75$, $H_a:\bar x>75$.
\QuestionsTasks[sec:general-worked]{{2}/{sol:notation}/{Correct the draft, formulate return hypotheses, and explain a revised cost claim of a reduction below 75 USD.}}
\subsubsection*{Worked solution}\label{sol:notation}
For population mean cost $\Pop{\mu_c}$, use $H_0:\mu_c=75$, $H_a:\mu_c>75$, right-tailed. $\Obs{\bar x}$ is evidence, not the parameter. For population mean return $\mu_r$ in percentage points, use $H_0:\mu_r=4$, $H_a:\mu_r\ne4$, two-tailed. A cost reduction changes $H_a$ to $\mu_c<75$, left-tailed, retaining the benchmark. Sampling and distribution conditions matter for later testing, but do not select direction.
\secInicio[prob:notation][sec:general-worked]
\section{Advanced Problems}\label{sec:advanced}
\SectionProblemsTableOfContents{\SectionProblemLink{prob:audit}{Audit a direction selected after sampling}}
\subsection{Audit a direction selected after sampling}\label{prob:audit}
A team planned to ask whether mean monthly household saving differs from 90 USD. After observing $\bar x=96$, it changes the alternative to an increase for a smaller p-value. Separately, a risk officer says the default rate is at most 5\% and asks whether that guarantee is violated.
\QuestionsTasks[sec:advanced]{{2}/{sol:audit}/{Formulate both legitimate questions and justify the tails.}}
\subsubsection*{Worked solution}\label{sol:audit}
For population mean saving $\Pop{\mu}$ retain $H_0:\mu=90$, $H_a:\mu\ne90$, two-tailed. Choosing $>$ after observing 96 makes direction depend on evidence. A future preregistered study could investigate an increase with new data. For population one-year default proportion $p$, a violation is $H_a:p>0.05$, right-tailed, with composite null $p\le0.05$ and calculation boundary $p=0.05$. The phrase at most describes the null-side guarantee, not a left-tailed alternative. Wording alone supplies no decision.
\secInicio[prob:audit][sec:advanced]
\end{document}

